% Chapter 1 — where it started: IMKG (ESWC 2023).

\gmenter{before}
\note{Chapter 1. Before anything I built: what existed when I arrived. Three stops: what a
knowledge graph is, what IMKG was, and what I set out to change.}

\gmgo{kg-picture}{A knowledge graph}
\begin{frame}{A knowledge graph, in one picture}
  \centering
  \begin{tikzpicture}[x=1cm, y=1cm,
      pic/.style={circle, line width=2pt, inner sep=0pt},
      ico/.style={circle, minimum size=1.1cm, inner sep=0pt, font=\fontsize{17}{17}\selectfont, text=white},
      nm/.style={font=\normalsize\bfseries, text=gmInk, inner sep=2pt, align=center},
      fact/.style={-{Stealth[length=2.6mm]}, line width=1.2pt, draw=gmBase, shorten <=1mm, shorten >=1mm},
      rel/.style={font=\small\itshape, text=gmInk2, fill=gmSurface, inner sep=1.2pt}]
    \node[pic, draw=kFrame, minimum size=2.1cm,
          path picture={\node at (path picture bounding box.center)
                        {\includegraphics[height=2.2cm]{figures/img/onedoes}};}] (meme) at (0,0) {};
    \node[nm, below=1mm of meme] {One Does Not Simply\\Walk into Mordor};
    \node[pic, draw=kWikidata, minimum size=1.25cm,
          path picture={\node at ($(path picture bounding box.center)+(.13,-.48)$)
                        {\includegraphics[height=2.95cm]{figures/img/onedoes}};}] (boromir) at (-3.8,1.35) {};
    \node[nm, left=1mm of boromir] {Boromir};
    \node[ico, fill=kWikidata] (sauron) at (-3.8,-1.45) {\faEye};
    \node[nm, left=1mm of sauron] {Sauron};
    \node[ico, fill=kWikidata] (mordor) at (3.8,1.35) {\faMountain};
    \node[nm, right=1mm of mordor] {Mordor};
    \node[ico, fill=kWikidata] (lotr) at (3.8,-1.45) {\faRing};
    \node[nm, below=1mm of lotr] {The Lord of the Rings};
    \node[ico, fill=kWikidata] (tolkien) at (6.7,-1.45) {\faFeather*};
    \node[nm, below=1mm of tolkien] {Tolkien};
    \node[ico, fill=kFrame, minimum size=.95cm, font=\fontsize{13}{13}\selectfont] (sub) at (0,2.45) {\faUsers};
    \node[nm, right=1mm of sub, align=left] {Lord of the Rings\\[-.6mm]{\small\mdseries\color{gmInk2}a Know Your Meme subculture}};
    \draw[fact] (meme) -- node[rel, sloped, above] {its picture shows} (boromir);
    \draw[fact] (meme) -- node[rel, sloped, above] {its text mentions} (mordor);
    \draw[fact] (boromir) -- node[rel, left] {enemy of} (sauron);
    \draw[fact] (mordor) -- node[rel, right] {appears in} (lotr);
    \draw[fact] (tolkien) -- node[rel, above] {wrote} (lotr);
    \draw[fact] (meme) -- node[rel, right] {part of} (sub);
  \end{tikzpicture}
  \vskip2mm
  {\large Things are \textbf{dots}. Facts are \textbf{arrows}. A computer can follow the arrows.}
  \gmsource{Redrawn from IMKG's own example (Tommasini, Ilievski, Wijesiriwardene, ESWC 2023, Fig.~7).}
  \note{
  \begin{itemize}
    \item The one idea everyone needs for the rest of the hour.
    \item A knowledge graph stores facts as arrows between things. ``The meme's picture shows
      Boromir.'' ``Its text mentions Mordor.'' ``Mordor appears in The Lord of the Rings.''
      ``Tolkien wrote it.''
    \item Because the arrows connect, a computer can answer questions nobody wrote down:
      ``which memes are about Tolkien's work?'' --- follow the arrows backwards.
    \item Blue dots are Know Your Meme entries, orange dots are Wikidata items: the two colours
      come back all talk.
    \item This picture is IMKG's own example --- the project I started from.
  \end{itemize}}
\end{frame}

\gmgo{imkg}{IMKG, 2023}
\begin{frame}{IMKG, the Internet Meme Knowledge Graph}
  \framesubtitle{ESWC 2023: a paper, its code, one data release --- four steps, run \textbf{once}}
  \centering
  \begin{tikzpicture}[x=1cm, y=1cm,
      stp/.style={circle, minimum size=1.3cm, fill=gm@col@before, text=white, inner sep=0pt,
                  font=\fontsize{17}{17}\selectfont},
      nm/.style={font=\large\bfseries, text=gmInk},
      what/.style={text width=3.0cm, align=center, font=\small, text=gmInk2},
      go/.style={-{Stealth[length=2.8mm]}, line width=1.2pt, draw=gmBase, shorten <=1.5mm, shorten >=1.5mm}]
    \foreach \i/\ic/\n/\w in {0/\faShapes/model/{media frames, memes, templates},
                              1/\faSpider/collect/{Know Your Meme in two days, imgflip, Wikidata},
                              2/\faMagic/enrich/{entities in the text and the picture},
                              3/\faProjectDiagram/integrate/{RDF; KYM and imgflip joined by hand}} {
      \node[stp] (s\i) at (\i*3.55,0) {\ic};
      \node[nm, below=1mm of s\i] (n\i) {\n};
      \node[what, below=0mm of n\i] {\w};}
    \foreach \i [count=\j] in {0,1,2} \draw[go] (s\i) -- (s\j);
  \end{tikzpicture}
  \vskip6mm
  \begin{columns}[t]
    \column{.24\textwidth}\centering\gmbig{\FImkgFrames}{KYM entries}
    \column{.24\textwidth}\centering\gmbig{\FImkgMemes}{imgflip memes}
    \column{.24\textwidth}\centering\gmbig{4.85M}{nodes}
    \column{.24\textwidth}\centering\gmbig{16.5M}{edges}
  \end{columns}
  \note{
  \begin{itemize}
    \item IMKG: the first knowledge graph of Internet memes, published at ESWC 2023 by Riccardo
      Tommasini (here at LIRIS), Filip Ilievski and Thilini Wijesiriwardene. Three things
      existed: the paper, the code on GitHub, and one data release on Zenodo (v0.9).
    \item Its four steps (Fig.~4 of the paper). Model: media frames, memes, templates. Collect: a
      Selenium crawler over all of Know Your Meme, about two days; imgflip's memes; Wikidata's
      556 memes as seeds. Enrich: entities from the text (DBpedia Spotlight) and the image
      (Google Vision), background facts from Wikidata (KGTK). Integrate: RML mappings to RDF;
      KYM and Wikidata joined on 276 memes, KYM and imgflip joined by hand and by title.
    \item Sources: Know Your Meme (the encyclopedia of memes), imgflip (the meme generator),
      Wikidata. Most of the size is imgflip: 15 of the 16.5 million edges are imgflip memes and
      their captions. The KYM part is 167,662 nodes and 914,941 edges.
    \item Then it was published. Nothing re-ran it.
  \end{itemize}}
\end{frame}

\gmgo{goals}{What I set out to build}
\begin{frame}{From a snapshot to a living graph}
  \begin{columns}[t]
    \column{.46\textwidth}
      {\large\bfseries\color{gmMuted}What I found}\par\vskip3mm
      \begin{tikzpicture}[x=1cm, y=-1.05cm,
          ic/.style={circle, minimum size=.85cm, fill=gmMuted!18!gmSurface, text=gmInk2,
                     font=\fontsize{12}{12}\selectfont, inner sep=0pt},
          tx/.style={anchor=west, text width=5.3cm, font=\normalsize, text=gmInk, align=left}]
        \node[ic] at (0,0) {\faCamera};    \node[tx] at (.6,0) {\textbf{A snapshot}\\[-.4mm]{\small\color{gmInk2}the site changes every day}};
        \node[ic] at (0,1) {\faEyeSlash};  \node[tx] at (.6,1) {\textbf{Most of each page unread}\\[-.4mm]{\small\color{gmInk2}links, references, the story}};
        \node[ic] at (0,2) {\faHandPaper}; \node[tx] at (.6,2) {\textbf{Bridges built by hand}\\[-.4mm]{\small\color{gmInk2}\FImkgPsixseven\ memes, \FImkgLinkedFrames\ entries}};
        \node[ic] at (0,3) {\faQuestion};  \node[tx] at (.6,3) {\textbf{No evidence on a link}\\[-.4mm]{\small\color{gmInk2}no score, no reason}};
      \end{tikzpicture}
    \column{.06\textwidth}\centering
      \vskip16mm{\fontsize{24}{24}\selectfont\color{gmBase}\faArrowRight}
    \column{.44\textwidth}
      {\large\bfseries\color{gmInk}What I set out to build}\par\vskip3mm
      \begin{tikzpicture}[x=1cm, y=-1.4cm,
          ic/.style={circle, minimum size=.85cm, fill=kFrame, text=white,
                     font=\fontsize{12}{12}\selectfont, inner sep=0pt},
          tx/.style={anchor=west, text width=5.1cm, font=\normalsize, text=gmInk, align=left}]
        \node[ic] at (0,0) {\faSync};      \node[tx] at (.6,0) {\textbf{A pipeline}, every month\\[-.4mm]{\small\color{gmInk2}redoing only what changed}};
        \node[ic] at (0,1) {\faSearch};    \node[tx] at (.6,1) {\textbf{Every page in depth}\\[-.4mm]{\small\color{gmInk2}and how sure we are of each fact}};
        \node[ic] at (0,2) {\faHandshake}; \node[tx] at (.6,2) {\textbf{IMKG's words}, kept\\[-.4mm]{\small\color{gmInk2}so its queries still run}};
      \end{tikzpicture}
  \end{columns}
  \vskip3mm
  {\normalsize\color{gmInk2}\itshape ``we plan recurrent releases of the KGs [\dots] to handle scalability in terms of
  volume and velocity''\upshape\enspace --- IMKG, Conclusion}
  \note{
  \begin{itemize}
    \item Be fair: IMKG was a research prototype and did what a paper needs. These are the
      things a \emph{living} graph needs, which is what I was asked to build.
    \item A snapshot: KYM adds and edits entries all the time; a one-off crawl is out of date
      the next week.
    \item Unread: IMKG kept About, Origin and Spread as text and annotated it; every link, every
      reference, every event in the story was left on the page.
    \item Hand-made bridges do not scale to new entries; and a link with no score or reason
      cannot be checked.
    \item Three goals, which come back in every chapter: a pipeline (re-runnable, incremental,
      the same inputs give the same graph); depth, with confidence on every derived fact; and
      compatibility --- MemeAtlas \emph{extends} IMKG, it does not replace it.
    \item And the authors said it themselves: recurrent releases, volume and velocity. That is
      the brief I took. Now back to the dog, and the first step: getting the pages.
  \end{itemize}}
\end{frame}
